%==========================================================================
%  TRES0312 Fysikk — Foiler-mal
%
%  MØrk/lys modus: kommenter ut én av de to linjene nedenfor.
%==========================================================================
\newif\ifdarkmode
\darkmodetrue      % ← mørk modus
%\darkmodefalse   % ← lys modus  (fjern % foran og legg % foran linjen over)

\documentclass[aspectratio=169, 12pt]{beamer}

\usepackage{amd-beamer}
\usepackage[european]{circuitikz}
\usetikzlibrary{overlay-beamer-styles}
\usetikzlibrary{calc}

%--------------------------------------------------------------------------
%  DOKUMENTINFO
%--------------------------------------------------------------------------
\title{Ladning, strøm og spenning}
\subtitle{TRES0312 --- kapittel 9.1 og 9.2}
\author{Ben David Normann}
\date{\today}

%--------------------------------------------------------------------------
%  SEKSJONSKONTEKST
%--------------------------------------------------------------------------
\setcounter{section}{9}   % Kapittel 9: "Elektrisitet"

\begin{document}

%--------------------------------------------------------------------------
%  SLIDE 1 — TITTELSIDE
%--------------------------------------------------------------------------
\begin{frame}[plain]
  \titlepage
  \dekkerdelkap{kapittel 9.1 og 9.2}
  \vspace{0.3cm}
  \begin{center}
  {\bfseries\color{repcol} Første time: elektrisk ladning og strøm. Etter pausen: spenning og kretser.}
  \end{center}
\end{frame}

%--------------------------------------------------------------------------
%  SLIDE 2 — Tre kjappe fra forrige gang
%--------------------------------------------------------------------------
\begin{frame}{}

  \repetisjon{Tre kjappe fra forrige gang}{%
    \begin{enumerate}[1.]
      \item Hvorfor må vi regne med temperaturen i kelvin i $pV = Nk_\text{B}T$?
      \item Hvor mye energi trengs for å varme $1.0\,\mathrm{kg}$ vann med $1\,^\circ$C?
      \item Hva er effekt? Hvilken enhet har effekt?
    \end{enumerate}
  }

\end{frame}

%==========================================================================
%  DEL 1: ELEKTRISK LADNING OG STRØM (9.1 OG 9.2)
%==========================================================================

%--------------------------------------------------------------------------
%  SLIDE 3 — Til ettertanke: ballongen
%--------------------------------------------------------------------------
\begin{frame}{Elektrisk ladning}

  \tanke{}{%
  Du gnir en ballong mot genseren din, og etterpå kan ballongen henge fast i taket. Hva er det som har skjedd?
  }
  \pause
  \svar{
  Elektroner er flyttet fra genseren til ballongen. Ballongen har fått et \textit{overskudd} av elektroner og er negativt ladet. Genseren har fått et \textit{underskudd} og er positivt ladet.
  }

\end{frame}

%--------------------------------------------------------------------------
%  SLIDE 4 — Ladning
%--------------------------------------------------------------------------
\begin{frame}{Elektrisk ladning}

  \begin{itemize}
    \item<1-> Atomer består av positive \textbf{protoner}, nøytrale nøytroner og negative \textbf{elektroner}.
    \item<2-> Et atom er \textbf{nøytralt}: like mange protoner som elektroner.
    \item<3-> Ladning er \textbf{bevart}. Den kan flyttes, men ikke lages eller ødelegges.
    \item<4-> Ladning måles i coulomb (C). Den minste frie ladningen er \textbf{elementærladningen}:
  \end{itemize}
  \onslide<4->{\graabox{$\displaystyle e = 1.602\cdot 10^{-19}\,\mathrm{C} \qquad\qquad q = N\cdot e$}}

\end{frame}

%--------------------------------------------------------------------------
%  SLIDE 5 — Aktivitet: ballongen
%--------------------------------------------------------------------------
%  Fasit: a) N = |q|/e = 0.20e-6/1.602e-19 = 1.25e12 ≈ 1.2e12 elektroner.
%         b) Genseren får +0.20 µC (ladningsbevaring).
\begin{frame}{Aktivitet}

  \aktivitet{Ladningen på en ballong}{%
  En ballong er gnidd mot en genser og har fått ladningen $q = -0.20\,\mathrm{\mu C}$.
  \begin{enumerate}[a)]
    \item Hvor mange elektroner har ballongen i overskudd?
    \item Hvilken ladning har genseren fått?
  \end{enumerate}
  }

\end{frame}

%--------------------------------------------------------------------------
%  SLIDE 6 — Ledere og isolatorer + orientering om kraft mellom ladninger
%--------------------------------------------------------------------------
\begin{frame}{Ledere og isolatorer}

  \begin{itemize}
    \item<1-> \textbf{Ledere} (metaller som kobber) har frie elektroner som kan bevege seg.
    \item<2-> \textbf{Isolatorer} (plast, gummi, glass) holder fast på elektronene sine.
    \item<3-> \textbf{Halvledere} (silisium) ligger imellom. De er grunnlaget for all elektronikk.
  \end{itemize}
  \onslide<4->{\merk{Til orientering (ikke pensum)}{%
  Likt fortegn frastøter, motsatt fortegn tiltrekker, og kraften avtar med avstanden. Derfor henger ballongen i taket: Den skyver elektronene i taket litt unna, og den positive overflaten som blir igjen, trekker på ballongen.
  }}

\end{frame}

%--------------------------------------------------------------------------
%  SLIDE 7 — Til ettertanke: lysbryteren
%--------------------------------------------------------------------------
\begin{frame}{Elektrisk strøm}

  \tanke{}{%
  Når du slår på lysbryteren, lyser lampa nesten med en gang. Farer elektronene fra bryteren til lampa på et øyeblikk?
  }
  \pause
  \svar{
  Nei, elektronene beveger seg typisk under en millimeter per sekund. Men ledningen er full av frie elektroner fra før. Alle begynner å bevege seg nesten samtidig, som vann i en hageslange som allerede er full.
  }

\end{frame}

%--------------------------------------------------------------------------
%  SLIDE 8 — Definisjon: strøm
%--------------------------------------------------------------------------
\begin{frame}{Elektrisk strøm}

  \defnat{9.2.1}{Elektrisk strøm}{
  Når ladningen $q$ passerer et tverrsnitt i løpet av tiden $t$, er strømmen
  \[
    I = \frac{q}{t} \qquad [I] = \mathrm{A} = \mathrm{C/s}
  \]
  }
  \pause
  \merk{Strømretning}{%
  Vi sier at strømmen går fra $+$ til $-$ utenfor batteriet. Elektronene beveger seg faktisk motsatt vei.
  }

\end{frame}

%--------------------------------------------------------------------------
%  SLIDE 9 — Aktivitet: lommelykt
%--------------------------------------------------------------------------
%  Fasit: a) t = 10 min = 600 s. q = 0.50 * 600 = 3.0e2 C.
%         b) N = 300/1.602e-19 = 1.87e21 ≈ 1.9e21 elektroner.
%         c) I = q/t = 1.0/2.0e-3 = 5.0e2 A.
\begin{frame}{Aktivitet}

  \aktivitet{Strøm og ladning}{%
  \begin{enumerate}[a)]
    \item En lommelykt trekker strømmen $0.50\,\mathrm{A}$ i $10\,\mathrm{min}$. Hvor mye ladning har gått gjennom lyspæra?
    \item Hvor mange elektroner tilsvarer det?
    \item En gnist overfører ladningen $1.0\,\mathrm{C}$ på $2.0\,\mathrm{ms}$. Hvor stor er strømmen?
  \end{enumerate}
  }

\end{frame}

%--------------------------------------------------------------------------
%  SLIDE 10 — Oppsummering: del 1
%--------------------------------------------------------------------------
\begin{frame}{}

  \oppsum{Elektrisk ladning og strøm}{%
  \begin{enumerate}[1.]
    \item[] Ladning er bevart. Et legeme blir ladet ved å få overskudd eller underskudd av elektroner.
    \item[] $q = N\cdot e$, med $e = 1.602\cdot 10^{-19}\,\mathrm{C}$.
    \item[] Strøm er ladning per tid: $I = q/t$, målt i ampere.
  \end{enumerate}
  }

\end{frame}

%--------------------------------------------------------------------------
%  SLIDE 11 — PAUSE
%--------------------------------------------------------------------------
\begin{frame}
\begin{center}
  \Huge{Pause}\\[8pt]
  {\large Etter pausen: spenning og kretser}
\end{center}
\end{frame}

%==========================================================================
%  DEL 2: SPENNING OG ELEKTRISKE KRETSER (9.2)
%==========================================================================

%--------------------------------------------------------------------------
%  SLIDE 13 — Definisjon: spenning
%--------------------------------------------------------------------------
\begin{frame}{Spenning}

  \defnat{9.2.2}{Spenning}{
  Spenningen $U$ mellom to punkter er arbeidet per ladning som trengs for å flytte ladningen fra det ene punktet til det andre:
  \[
    U = \frac{W}{q} \qquad [U] = \mathrm{V} = \mathrm{J/C}
  \]
  }
  \pause
  Tenk på spenning som \textbf{høydeforskjell}: Batteriet løfter ladningene opp, og strømmen «renner nedover» gjennom kretsen. Et batteri på $9.0\,\mathrm{V}$ gir $9.0\,\mathrm{J}$ til hver coulomb.

\end{frame}

%--------------------------------------------------------------------------
%  SLIDE 14 — Aktivitet: ladekabel
%--------------------------------------------------------------------------
%  Fasit: a) t = 45 min = 2700 s. q = 2.0 * 2700 = 5.4e3 C.
%         b) N = 5.4e3/1.602e-19 = 3.4e22.
%         c) W = U q = 5.0 * 5.4e3 = 2.7e4 J.
\begin{frame}{Aktivitet}

  \aktivitet{Strøm i en ladekabel}{%
  En mobil lades med strømmen $2.0\,\mathrm{A}$ i $45\,\mathrm{min}$ fra en lader på $5.0\,\mathrm{V}$.
  \begin{enumerate}[a)]
    \item Hvor mye ladning har gått gjennom kabelen?
    \item Hvor mange elektroner tilsvarer det?
    \item Hvor mye energi har laderen levert?
  \end{enumerate}
  }

\end{frame}

%--------------------------------------------------------------------------
%  SLIDE 15 — Kretser og symboler
%--------------------------------------------------------------------------
\begin{frame}{Elektriske kretser}

  \begin{columns}[c]
    \begin{column}{0.50\textwidth}
      \begin{itemize}
        \item<1-> Det går bare strøm i en \textbf{lukket} krets.
        \item<2-> \textbf{Amperemeter} kobles \textit{i serie}: strømmen må gå gjennom det.
        \item<3-> \textbf{Voltmeter} kobles \textit{i parallell} over komponenten.
      \end{itemize}
    \end{column}
    \begin{column}{0.48\textwidth}
      \centering
      \begin{circuitikz}[color=textcolor, scale=0.9, transform shape]
        \draw (0,0) to[battery1, l=$U$] (0,3)
              to[ammeter, l=$A$] (3,3)
              to[lamp] (3,0) -- (0,0);
        \draw (3,3) -- (4.6,3) to[voltmeter, l=$V$] (4.6,0) -- (3,0);
      \end{circuitikz}
    \end{column}
  \end{columns}

\end{frame}

%--------------------------------------------------------------------------
%  SLIDE 16 — Til ettertanke: amperemeter og voltmeter
%--------------------------------------------------------------------------
\begin{frame}{Elektriske kretser}

  \tanke{}{%
  Et godt amperemeter har svært liten resistans, mens et godt voltmeter har svært stor resistans. Hvorfor?
  }
  \pause
  \svar{
  Måleinstrumentet skal ikke endre kretsen. Et amperemeter i serie må slippe strømmen lett gjennom. Et voltmeter i parallell skal ikke «stjele» strøm fra komponenten det måler over.
  }

\end{frame}

%--------------------------------------------------------------------------
%  SLIDE 17 — Oppsummering: del 2
%--------------------------------------------------------------------------
\begin{frame}{}

  \oppsum{Spenning og kretser}{%
  \begin{enumerate}[1.]
    \item[] Strøm er ladning per tid: $I = q/t$.
    \item[] Spenning er energi per ladning: $U = W/q$.
    \item[] Amperemeter i serie, voltmeter i parallell.
  \end{enumerate}
  }

\end{frame}

%--------------------------------------------------------------------------
%  SLIDE 18 — Neste gang
%--------------------------------------------------------------------------
\begin{frame}
\begin{center}
  \Huge{Neste gang: Ohms lov, koblinger og effekt}
  \vspace{6mm}
\end{center}
\end{frame}

\end{document}
